\textbf{COLUNI 2026 -- Questão 13}\par\medskip Numa Gincana de Desafios Matemáticos, Cirbélia e Virgínia receberam, cada uma delas, um cartão personalizado contendo uma expressão numérica, conforme os modelos abaixo.

\begin{center}
\fbox{\begin{minipage}[c][3.5cm][c]{0.42\linewidth}\centering
$\dfrac{3^1-(-1)^2+2^2}{2^{-1}+2^{-2}-2^{-3}}=$\\[1.2em]
Cirbélia
\end{minipage}}\hfill
\fbox{\begin{minipage}[c][3.5cm][c]{0.42\linewidth}\centering
$\dfrac{1+4^{\frac12}-\sqrt[3]{-8}}{1+\sqrt{0,04}+\sqrt[3]{-1}}=$\\[1.2em]
Virgínia
\end{minipage}}
\end{center}

Cirbélia calculou o valor da expressão numérica contida em seu cartão e escreveu $\dfrac{19}{2}$, depois do sinal de igualdade. Virgínia também calculou o valor da expressão numérica contida em seu cartão e escreveu 25, após o sinal de igualdade.

É CORRETO afirmar que:
\begin{enumerate}
\item ambas erraram.
\item ambas acertaram.
\item apenas Cirbélia acertou.
\item apenas Virgínia acertou.
\end{enumerate}
